\chapter{Derived Categories of CY Manifolds}
\label{ch:derived-cy}

Homological mirror symmetry identifies $\Fuk(X) \simeq D^b(\Coh(X^\vee))$.  Under the functor $\Phi$, this becomes an equivalence $A_{\Fuk(X)} \simeq A_{D^b(\Coh(X^\vee))}$ of quantum chiral algebras.  This chapter develops the algebraic side: the CY structure on $D^b(\Coh(X))$, the chiral algebra extracted by $\Phi$, exceptional collections, and Bridgeland stability conditions.

\section{$D^b(\Coh(X))$ for CY manifolds}
\label{sec:db-coh-cy}

For a smooth projective CY manifold $X$ of dimension $d$, the bounded derived category $D^b(\Coh(X))$ is CY of dimension $d$ with CY trace induced by Serre duality.  The Hochschild homology $\HH_\bullet(D^b(\Coh(X))) \simeq H^*(X, \Omega^\bullet_X)$ (HKR isomorphism) recovers the Hodge cohomology, and the CY trace is the Grothendieck residue.  The functor $\Phi$ applied to $D^b(\Coh(X))$ produces the B-model chiral algebra, with $\kappa = \chi(\cO_X)$ for CY surfaces.

\section{Homological mirror symmetry and the chiral algebra}
\label{sec:hms-chiral}

The homological mirror symmetry conjecture (Kontsevich) identifies $\Fuk(X) \simeq D^b(\Coh(X^\vee))$.  Under the CY-to-chiral functor $\Phi$, this becomes an equivalence of quantum chiral algebras:
\[
  A_{\Fuk(X)} \simeq A_{D^b(\Coh(X^\vee))}.
\]

\section{Exceptional collections and chiral algebras}
\label{sec:exceptional-chiral}

When $D^b(\Coh(X))$ admits a full exceptional collection $\langle E_1, \ldots, E_n \rangle$, the CY-to-chiral functor reduces to a finite-dimensional computation: the $\Ainf$-algebra of the collection determines a quiver with potential, and $\Phi$ produces the critical CoHA of the quiver (connecting to Chapter~\ref{ch:toric-coha}).  This provides the simplest non-trivial test cases for Theorem~CY-A.

\section{Stability conditions and the chiral moduli}
\label{sec:stability-chiral}

Bridgeland stability conditions on $D^b(\Coh(X))$ parametrise the space of t-structures, providing a moduli space $\Stab(X)$ over which the chiral algebra $\Phi(D^b(\Coh(X)))$ varies.  Wall-crossing in $\Stab(X)$ corresponds to mutations of the bar complex; the global chiral algebra is assembled as a homotopy colimit over stability chambers (Programme~A of the $K3 \times E$ programme, Chapter~\ref{ch:toroidal-elliptic}).
